% Part: computability
% Chapter: recursive-functions
% Section: general-recursive-functions

\documentclass[../../../include/open-logic-section]{subfiles}

\begin{document}

\olfileid[hi]{cmp}{rec}{gen}
\olsection{सामान्य पुनरावर्ती फलन}

पूर्ण फलनों का समुच्चय पाने का एक अन्य ढंग है। किसी पूर्ण फलन
$f(x,\vec z)$ को \emph{नियमित} कहें यदि प्राकृतिक संख्याओं के प्रत्येक
अनुक्रम $\vec z$ के लिए कोई $x$ ऐसा हो कि $f(x,\vec z) = 0$। दूसरे
शब्दों में, नियमित फलन ठीक वे फलन हैं जिन पर अपरिबद्ध खोज लागू करके अंत
में पूर्ण फलन मिलता है। हम सावधानीपूर्वक अपरिबद्ध खोज को नियमित फलनों तक
सीमित कर सकते हैं:

\begin{defn}
\ollabel{defn:general-recursive}
\emph{सामान्य पुनरावर्ती फलनों} का समुच्चय प्राकृतिक संख्याओं से
प्राकृतिक संख्याओं में (विभिन्न स्थान-संख्याओं वाले) फलनों का वह लघुतम
समुच्चय है जिसमें शून्य, उत्तराधिकारी और प्रक्षेपण आते हैं तथा जो संयोजन,
आदिम पुनरावर्तन और \emph{नियमित} फलनों पर लागू अपरिबद्ध खोज के अधीन
संवृत है।
\end{defn}

स्पष्टतः प्रत्येक सामान्य पुनरावर्ती फलन पूर्ण है।
\olref{defn:general-recursive} और \olref[par]{defn:recursive-fn} के बीच
अंतर यह है कि बाद वाले में प्रक्रिया के दौरान आंशिक पुनरावर्ती फलनों का
उपयोग करने की अनुमति है; एकमात्र अपेक्षा यह है कि अंत में मिलने वाला फलन
पूर्ण हो। इसलिए ऐतिहासिक अवशेष ``सामान्य'' शब्द अनुपयुक्त नाम है; ऊपर से
देखने पर \olref{defn:general-recursive},
\olref[par]{defn:recursive-fn} से \emph{कम} सामान्य है। किंतु सौभाग्य से
यह अंतर आभासी है; यद्यपि परिभाषाएँ भिन्न हैं, सामान्य पुनरावर्ती फलनों का
समुच्चय और पुनरावर्ती फलनों का समुच्चय एक ही है।

\end{document}

